% Options for packages loaded elsewhere
\PassOptionsToPackage{unicode}{hyperref}
\PassOptionsToPackage{hyphens}{url}
\documentclass[
  english,
]{article}
\usepackage{xcolor}
\usepackage{amsmath,amssymb}
\setcounter{secnumdepth}{-\maxdimen} % remove section numbering
\usepackage{iftex}
\ifPDFTeX
  \usepackage[T1]{fontenc}
  \usepackage[utf8]{inputenc}
  \usepackage{textcomp} % provide euro and other symbols
\else % if luatex or xetex
  \usepackage{unicode-math} % this also loads fontspec
  \defaultfontfeatures{Scale=MatchLowercase}
  \defaultfontfeatures[\rmfamily]{Ligatures=TeX,Scale=1}
\fi
\usepackage{lmodern}
\ifPDFTeX\else
  % xetex/luatex font selection
\fi
% Use upquote if available, for straight quotes in verbatim environments
\IfFileExists{upquote.sty}{\usepackage{upquote}}{}
\IfFileExists{microtype.sty}{% use microtype if available
  \usepackage[]{microtype}
  \UseMicrotypeSet[protrusion]{basicmath} % disable protrusion for tt fonts
}{}
\makeatletter
\@ifundefined{KOMAClassName}{% if non-KOMA class
  \IfFileExists{parskip.sty}{%
    \usepackage{parskip}
  }{% else
    \setlength{\parindent}{0pt}
    \setlength{\parskip}{6pt plus 2pt minus 1pt}}
}{% if KOMA class
  \KOMAoptions{parskip=half}}
\makeatother
\ifLuaTeX
\usepackage[bidi=basic,shorthands=off]{babel}
\else
\usepackage[bidi=default,shorthands=off]{babel}
\fi
\ifLuaTeX
  \usepackage{selnolig} % disable illegal ligatures
\fi
\setlength{\emergencystretch}{3em} % prevent overfull lines
\providecommand{\tightlist}{%
  \setlength{\itemsep}{0pt}\setlength{\parskip}{0pt}}
\usepackage{bookmark}
\IfFileExists{xurl.sty}{\usepackage{xurl}}{} % add URL line breaks if available
\urlstyle{same}
\hypersetup{
  pdftitle={Orthonormal bases and orthogonal projections},
  pdflang={en},
  hidelinks,
  pdfcreator={LaTeX via pandoc}}

\title{Orthonormal bases and orthogonal projections}
\author{}
\date{}

\begin{document}
\maketitle

\section{Orthonormal bases and orthogonal
projections}\label{orthonormal-bases-and-orthogonal-projections}

This lesson supplies the inner-product facts used in
\href{https://kokunoyumeto.github.io/open-mathematics-courses/courses/RT-FIN/representations-and-complete-reducibility.html\#2-averaging-separates-invariant-pieces}{Representations
and complete reducibility, Proposition 2.1 and its following paragraph}.
It belongs to the linear-algebra foundation. It does not require the
functional-analysis course.

\emph{Proofs, explanations, examples, solutions and programme
integration produced by OpenAI Codex --- GPT-6 Astra, Ultra effort,
October 2026. Self-checked by the writing AI; no human or independent
review is claimed. These are standard results, not discoveries.}

\subsection{Prerequisites and sources}\label{prerequisites-and-sources}

We use real and complex arithmetic, complex conjugation, positive square
roots of positive real numbers, finite induction, vector-space axioms
and finite sums. The programme's earlier lesson
\href{../basis-projection-bridge/\#extending-a-basis-and-choosing-a-complement}{From
bases to projections} proves that every subspace of a finite-dimensional
vector space has a finite basis and that a basis extends to the whole
space. No completeness, limit of vectors, choice of an infinite basis,
or projection theorem is assumed in the proofs of Theorems 1--2.
Exercise 3 is an optional comparison using square-summable sequences and
convergence of geometric series.

John M. Erdman's \emph{Functional Analysis and Operator Algebras: An
Introduction}, version 4 October 2015, Chapter 1, states Gram--Schmidt
orthonormalization and the orthogonal-decomposition corollary (source
label \texttt{0001514}). The English and Indonesian source passages are
part of
\href{https://kokunoyumeto.github.io/functional-analysis-erdman-id/}{the
programme's Erdman edition}. Those particular passages contain
statements without proofs. This bridge supplies the elementary proofs
and connects them to their use in representation theory. It uses
Erdman's convention, linear in the first variable. The
\href{https://web.pdx.edu/~erdman/FAOA/functional_analysis_operator_algebras_pdf.pdf}{author's
source edition} and the programme edition retain his authorship; he has
not reviewed or endorsed this bridge.

This exposition is offered under
\href{https://creativecommons.org/licenses/by-sa/4.0/}{Creative Commons
Attribution--ShareAlike 4.0}, preserving the terms of the Erdman
material it develops. The linked basis-extension lesson retains its own
CC BY-SA 2.5 licence; its text is not copied into this lesson. Exact
source identities and the bounded comparison with other programme proofs
are recorded in \url{SOURCE_REVIEW.json}.

\subsection{Inner products and
orthogonality}\label{inner-products-and-orthogonality}

Let \(\mathbb F\) be \(\mathbb R\) or \(\mathbb C\). An \textbf{inner
product} on an \(\mathbb F\)-vector space \(V\) is a map
\(\langle\ ,\ \rangle:V\times V\to\mathbb F\) such that

\[
\langle ax+by,z\rangle=a\langle x,z\rangle+b\langle y,z\rangle,
\qquad \langle y,x\rangle=\overline{\langle x,y\rangle},
\]

and \(\langle x,x\rangle\) is a positive real number for every nonzero
\(x\). Here \(a,b\in\mathbb F\); conjugation does nothing over
\(\mathbb R\). The two identities imply

\[
\langle x,ay+bz\rangle=\overline a\langle x,y\rangle+\overline b\langle x,z\rangle.
\]

They also give \(\langle0,x\rangle=\langle x,0\rangle=0\). Put
\(\|x\|=\sqrt{\langle x,x\rangle}\). In particular \(\|ax\|=|a|\|x\|\),
and \(\|x\|=0\) exactly when \(x=0\). We will need only these facts and
direct expansions, not the triangle inequality.

Vectors are \textbf{orthogonal} when their inner product is zero. For a
subspace \(U\), define

\[
U^\perp=\{v\in V:\langle v,u\rangle=0\text{ for every }u\in U\}.
\]

Linearity in the first variable makes \(U^\perp\) a subspace. Conjugate
symmetry makes orthogonality symmetric. A family \((e_j)\) is
\textbf{orthonormal} when \(\langle e_i,e_j\rangle\) is one for \(i=j\)
and zero otherwise. For a finite linear combination, taking its inner
product with \(e_i\) recovers its coefficient. Thus every orthonormal
family is linearly independent.

\subsection{Gram--Schmidt with the coefficients in the correct
order}\label{gramschmidt-with-the-coefficients-in-the-correct-order}

\textbf{Theorem 1.} Let \(v_1,v_2,\ldots\) be a linearly independent
finite or countably infinite sequence in a real or complex inner-product
space. There is an orthonormal sequence of the same length such that,
for every available positive integer \(j\),

\[
\operatorname{span}(e_1,\ldots,e_j)=\operatorname{span}(v_1,\ldots,v_j).
\]

No completeness or finite dimension of the ambient space is needed for
this statement. For an empty sequence, take the empty sequence.

\textbf{Proof.} Suppose the vectors through \(e_{j-1}\) have been
constructed, with the asserted span equality. Set

\[
w_j=v_j-\sum_{i=1}^{j-1}\langle v_j,e_i\rangle e_i.
\]

For \(k<j\), linearity and orthonormality give

\[
\langle w_j,e_k\rangle
=\langle v_j,e_k\rangle-\sum_{i=1}^{j-1}\langle v_j,e_i\rangle\langle e_i,e_k\rangle=0.
\]

If \(w_j=0\), then \(v_j\) belongs to the span of the previous \(e_i\),
hence of the previous \(v_i\). This contradicts independence. Therefore
\(\|w_j\|>0\), and we may define

\[
e_j=\frac{w_j}{\|w_j\|}.
\]

This vector has inner product one with itself and zero with each
previous vector. It belongs to the span of \(v_1,\ldots,v_j\).
Conversely,

\[
v_j=\|w_j\|e_j+\sum_{i=1}^{j-1}\langle v_j,e_i\rangle e_i
\]

belongs to the span of \(e_1,\ldots,e_j\). Together with the induction
hypothesis this proves both inclusions of spans. At \(j=1\) the sums are
empty, so the same argument starts the induction. In the countably
infinite case this construction defines each term by a finite
calculation; it makes no assertion about convergence of infinite linear
combinations. \(\square\)

\textbf{Corollary 1.} Every finite-dimensional real or complex
inner-product space has an orthonormal basis, and so does every subspace
of it.

\textbf{Proof.} Apply Theorem 1 to a finite basis, using the earlier
basis-extension lesson for a subspace's finite basis. Span equality
gives a spanning orthonormal family. In dimension zero the empty family
is the required basis. \(\square\)

\subsection{Orthogonal projection and
decomposition}\label{orthogonal-projection-and-decomposition}

\textbf{Theorem 2.} Let \(U\) be a finite-dimensional subspace of any
real or complex inner-product space \(V\). Then every \(v\in V\) has a
unique expression

\[
v=u+w,\qquad u\in U,\quad w\in U^\perp.
\]

Thus \(V=U\oplus U^\perp\). In particular this holds for every subspace
of a finite-dimensional \(V\). If \((e_1,\ldots,e_m)\) is any
orthonormal basis of \(U\), the orthogonal projection is

\[
Pv=\sum_{i=1}^m\langle v,e_i\rangle e_i.
\]

The map is independent of the chosen orthonormal basis. It is linear,
satisfies \(P^2=P\), has image \(U\) and kernel \(U^\perp\), and obeys

\[
\langle Px,y\rangle=\langle x,Py\rangle,
\qquad \|Pv\|^2+\|v-Pv\|^2=\|v\|^2.
\]

\textbf{Proof.} The restriction of the inner product to \(U\) is
positive definite, so Corollary 1 supplies its orthonormal basis. The
formula plainly takes values in \(U\). For every \(k\), expansion gives

\[
\langle v-Pv,e_k\rangle
=\langle v,e_k\rangle-\sum_i\langle v,e_i\rangle\langle e_i,e_k\rangle=0.
\]

Conjugate-linearity in the second variable then gives
\(v-Pv\in U^\perp\). If \(z\in U\cap U^\perp\), taking its inner product
with itself gives \(z=0\). Subtracting two decompositions therefore
proves uniqueness. This also proves independence from the basis, since
every basis formula produces a decomposition with the same two specified
subspaces.

Linearity of \(P\) follows from linearity of each coefficient in \(v\).
For \(u\in U\), the unique decomposition is \(u=u+0\), so \(Pu=u\).
Consequently the image is exactly \(U\) and \(P^2=P\). Its formula
vanishes on \(U^\perp\); conversely \(Pv=0\) in the decomposition
implies \(v\in U^\perp\). This identifies its kernel.

Orthogonality gives

\[
\langle Px,y\rangle=\langle Px,Py\rangle=\langle x,Py\rangle.
\]

Expanding the inner product of \(v=Pv+(v-Pv)\) with itself makes the two
cross terms zero and proves the squared-length identity. It gives
\(\|Pv\|\le\|v\|\) as well. No norm-completeness assertion was used. If
\(U=0\), the empty sum gives \(P=0\); if \(U=V\), then \(P=1_V\). These
statements include \(V=0\). \(\square\)

\textbf{Corollary 2 (nearest vector).} For every \(v\in V\), \(Pv\) is
the unique vector of \(U\) minimizing distance to \(v\).

\textbf{Proof.} For \(u\in U\), the vectors \(v-Pv\) and \(Pv-u\) are
orthogonal. Expanding their squared length gives

\[
\|v-u\|^2=\|v-Pv\|^2+\|Pv-u\|^2.
\]

The last summand is nonnegative and vanishes exactly for \(u=Pv\).
\(\square\)

\subsection{The two uses in representation
theory}\label{the-two-uses-in-representation-theory}

Let a finite group act linearly on a finite-dimensional complex vector
space. Choose any basis and define an initial inner product by
\((x,y)_0=\sum_j x_j\overline{y_j}\) in its unique coordinates. It is
linear first, Hermitian, and positive definite because a nonzero vector
has a nonzero coordinate. The representation-theory lesson constructs
the invariant average

\[
\langle x,y\rangle_G=\frac1{|G|}\sum_{g\in G}(gx,gy)_0.
\]

Each group element acts invertibly, so every term on the diagonal is
positive when \(x\ne0\). Right multiplication permutes the group
elements, which proves invariance. Corollary 1 supplies the orthonormal
basis used at the end of Proposition 2.1. If \(M\) is the matrix of a
group element in that basis, invariance gives \(M^*M=I\): the inner
products of its columns are those of the original basis vectors, with
complex conjugation on the second argument (equivalently their
conjugates give the usual matrix entries of \(M^*M\)).

For an invariant subspace \(U\), a vector \(v\in U^\perp\), an element
\(g\in G\), and \(u\in U\), invariance gives

\[
\langle gv,u\rangle_G=\langle v,g^{-1}u\rangle_G=0,
\]

since \(g^{-1}u\in U\). Thus \(U^\perp\) is invariant. Theorem 2 now
supplies the decomposition used immediately after Proposition 2.1. The
orthogonal projection also commutes with the action: applying \(g\) to
both parts of the unique decomposition preserves their subspaces, so
\(P(gv)=gPv\).

These arguments close these two particular linear-algebra uses. They do
not prove every other prerequisite of the representation-theory course.
Averaging over a general field has its own characteristic condition, and
the earlier basis/projection bridge supplies its different, purely
algebraic starting projection.

\subsection{A complex example}\label{a-complex-example}

In \(\mathbb C^2\) use
\(\langle x,y\rangle=x_1\overline{y_1}+x_2\overline{y_2}\). Starting
with \(v_1=(1,i)\), \(v_2=(0,1)\), Theorem 1 gives

\[
e_1=\frac{(1,i)}{\sqrt2},\quad
\langle v_2,e_1\rangle=-\frac{i}{\sqrt2},\quad
w_2=(i/2,1/2),\quad e_2=\frac{(i,1)}{\sqrt2}.
\]

The two vectors have squared length one and
\(\langle e_1,e_2\rangle=(-i+i)/2=0\). For \(U=\mathbb C(1,i)\), Theorem
2 gives

\[
P=\frac12\begin{pmatrix}1&-i\\i&1\end{pmatrix}.
\]

For example \(P(0,1)=(-i/2,1/2)\) and \((0,1)-P(0,1)=(i/2,1/2)\).
Reversing the arguments in the coefficient would not give this
decomposition. In complex spaces, the coefficient order is substantive.

\subsection{Exercises with solutions}\label{exercises-with-solutions}

\textbf{Exercise 1.} Explain why a positive semidefinite form is not
enough for Gram--Schmidt as stated.

\textbf{Solution.} On the one-dimensional space \(\mathbb C\), the zero
form is positive semidefinite. The singleton \((1)\) is linearly
independent but its squared length is zero, so it cannot be normalized
to length one. The proof needs positive definiteness precisely when it
divides by \(\|w_j\|\).

\textbf{Exercise 2.} For the complex matrix \(P\) above, verify directly
that \(P^2=P=P^*\), find its image and kernel, and decompose \((1,0)\).

\textbf{Solution.} Write \(A=2P\). Multiplication gives

\[
A^2=\begin{pmatrix}2&-2i\\2i&2\end{pmatrix}=2A,
\qquad A^*=A.
\]

Thus \(P^2=P=P^*\). Its image is spanned by \((1,i)\), since its first
column is half that vector and its second is \(-i\) times the first. Its
kernel is the line \(x-iy=0\), spanned by \((i,1)\). Finally

\[
(1,0)=(1/2,i/2)+(1/2,-i/2)
\]

has its first term in the image and its second in the kernel. Their
inner product is zero.

\textbf{Exercise 3.} Why does Theorem 2 keep the assumption that \(U\)
is finite-dimensional, even when the ambient space is allowed to be
infinite-dimensional?

\textbf{Solution.} In the usual sequence space \(\ell^2(\mathbb N)\),
let \(U\) consist of sequences with only finitely many nonzero entries.
If \(x\in U^\perp\), taking inner products with each coordinate vector
gives \(x_n=0\) for every \(n\), so \(U^\perp=0\). But the sequence
\((2^{-n})_{n\ge1}\) belongs to \(\ell^2\), because
\(\sum_{n\ge1}4^{-n}=1/3\), and does not belong to \(U\). Therefore
\(\ell^2\ne U\oplus U^\perp\). This does not contradict the later
projection theorem for \textbf{closed} subspaces of a Hilbert space;
this \(U\) is not closed, since its truncations converge to that same
sequence.

For the broader result, see
\href{https://kokunoyumeto.github.io/open-mathematics-courses/courses/foundations-of-von-neumann-algebras/hilbert-spaces-and-compact-operators.html\#OA-FND-HS-02}{Hilbert
spaces and compact operators, Theorems 2.1--2.2}. That proof uses
completeness and closedness to obtain a nearest vector. It remains the
programme's general Hilbert-space treatment; the elementary proof here
does not replace or narrow it.

\end{document}
